\input{../tex-inputs/latex-leseansicht-vorspann}

\section[Arthur Schnitzler an Gustav Schwarzkopf, 12. 8. 1913]{L04512 Arthur Schnitzler an Gustav Schwarzkopf, 12. 8. 1913}
\nopagebreak\mylabel{L04512v}
\rehead{ }\normalsize\beginnumbering\briefempfaengerindex{Schwarzkopf, Gustav@\textsc{Schwarzkopf, Gustav}!zzzSchnitzler, Arthur@\emph{von Arthur Schnitzler}!1913-08-121@{12. 8. 1913}|(be}
\toendnotes[C]{\smallbreak\pagebreak[2]}
\correspDesc{Versand  durch Arthur Schnitzler am 12. 8. 1913 in Brijuni
\newline{}Übermittlung  am 12. 8. 1913 in Brijuni
\newline{}Erhalt  durch Gustav Schwarzkopf im Zeitraum [13. 8. 1913 – 17. 8. 1913?] in Wien}\toendnotes[C]{\smallbreak}
\Standort{DLA, A:Schnitzler, HS.1985.1.1897.}
\physDesc{Bildpostkarte, 263 Zeichen
\newline{}Handschrift: Bleistift, deutsche Kurrent
\newline{}Versand: Stempel: »\nobreak{}\oindex{Brijuni@\textbf{Brijuni}|pwk}Brioni, 12. 8. 13\nobreak{}«.  }
\pstart
           \noindent{}\centering{}{\pb}\textcolor{gray}{\textbf{Brioni\oindex{Brijuni@\textbf{Brijuni}|pw}
                  :}}\pend
           
\pstart
           \centering{}\textcolor{gray}{\textbf{Venustempel in Val-Catena\oindex{Tempel der Venus@\textbf{Tempel der Venus}|pw}.}}\pend
           \pstart{}{\pb}\textsc{Herrn Gustav Schwarzkopf}\pend{}\pstart{}Wien I\oindex{I., Innere Stadt@\textbf{I., Innere Stadt}|pw}\pend{}\pstart{}\textsc{Tiefer Graben 17}\oindex{Wien@\textbf{Wien}!I., Innere Stadt@\textbf{I., Innere Stadt}!Tiefer Graben 17@\textbf{Tiefer Graben 17}|pw}\pend{}{\bigskip}\vspace{1em}
\pstart
           \raggedleft{}\textsc{Brioni\oindex{Brijuni@\textbf{Brijuni}|pw}}, 12. 8. 913\pend
           \vspace{0.5em}
\pstart
           lieber Guſtav, laſſen Sie doch etwas
      von ſich hören! Wir dürften noch
      bis zum 22. 8{[}.{]} bleiben, u haben
               nur Engadiner\oindex{Engadin@\textbf{Engadin}|pw} Pläne für 
      den Fall{ }ſchönen Wetters.\pend
           \pstart Seien Sie von uns allen
      herzlichſt gegrüßt! Ihr \spacefill\mbox{Arthur}\pend{}\selectlanguage{ngerman}\endnumbering\briefempfaengerindex{Schwarzkopf, Gustav@\textsc{Schwarzkopf, Gustav}!zzzSchnitzler, Arthur@\emph{von Arthur Schnitzler}!1913-08-121@{12. 8. 1913}|)be}\mylabel{L04512h}
\begin{anhang}
\end{anhang}\newcommand{\dateiname}{L04512}\newcommand{\titel}{Arthur Schnitzler an Gustav Schwarzkopf, 12. 8. 1913}\newcommand{\editorInnen}{Herausgegeben von Jahnke, SelmaMüller, Martin Anton}\input{../tex-inputs/latex-leseansicht-abspann}
